problem:  
Let \((M,\Delta,g)\) be a compact, connected, bracket-generating **contact sub-Riemannian manifold** of corank \(1\), with the contact distribution \(\Delta = \ker \alpha\) arising from a globally defined contact form \(\alpha\). Let \(H(x,p) = \tfrac{1}{2}\langle p|_{\Delta_x}, p|_{\Delta_x} \rangle_g\) be the associated sub-Riemannian Hamiltonian on \(T^*M\), and let \(\phi_t^H\) denote its Hamiltonian flow.  

Assume that there exists a constant \(T>0\) such that **for every unit covector \(p \in T^*_x M\)**, the normal geodesic \(\gamma_p(t) = \pi(\phi_t^H(p))\) is closed with period \(T\). Equivalently, the Hamiltonian flow \(\phi_t^H\) on the unit (“energy 1/2”) cotangent bundle is **periodic of period \(T\)**.  

Conjectural question:  
Must \((M,\Delta,g)\) necessarily arise from a **regular Sasakian manifold**, i.e., is there a choice of contact form \(\tilde{\alpha}\) co-oriented with \(\alpha\) and a compatible metric \(g_{\tilde{\alpha}}\) such that the Reeb flow of \(\tilde{\alpha}\) generates the same periodic geodesic dynamics, and \(M\) is a **Boothby–Wang circle bundle over a Kähler manifold** with constant holomorphic sectional curvature?  

Why is it a "good" problem:  
This problem formulates a *sub-Riemannian analogue of the classical Zoll rigidity conjecture* in Riemannian geometry, but in the profoundly more intricate contact–Hamiltonian setting. It explicitly sharpens the link between periodic sub-Riemannian Hamiltonian flows and the existence of Sasakian (and hence Kähler) structures, extending the well-known Boothby–Wang correspondence into the analytical domain of nonintegrable metrics.  

The problem embodies deep cross-disciplinary connections:  
- **Hamiltonian dynamics and geometric control theory:** Periodicity of the sub-Riemannian geodesic flow imposes nontrivial integrability constraints.  
- **Contact and symplectic geometry:** The assumption of periodic flow evokes questions about the regularity of the Reeb dynamics and the topology of the contact manifold.  
- **Complex and metric geometry:** The conjectured Sasakian rigidity ties the contact structure to Kähler geometry on the quotient space.  

Simplicity of statement, depth of consequences, and the potential to expose a new “Zoll-type rigidity” phenomenon in sub-Riemannian geometry make the question both open and conceptually deep—a narrow entrance leading to a landscape that integrates contact topology, Hamiltonian mechanics, and complex differential geometry.